\documentclass[a4paper]{article}
% Español (México) con XeLaTeX: fontspec para fuentes Unicode, polyglossia
% para la división silábica y los nombres del documento en la variante mexicana.
\usepackage{fontspec}
\usepackage{polyglossia}
\setdefaultlanguage[variant=mexican]{spanish}
% Los nombres chinos van como «pinyin (original)»: xeCJK da a los caracteres
% Han su propia fuente; sin ella XeLaTeX los omite con un aviso.
% FandolSong no cubre algunos caracteres de nombres (U+73FA); WenQuanYi Zen Hei sí.
\usepackage[AutoFallBack=true]{xeCJK}
\setCJKmainfont{FandolSong-Regular.otf}
\setCJKfallbackfamilyfont{\CJKrmdefault}{WenQuanYi Zen Hei}
% polyglossia no traduce los nombres de listings.
\AtBeginDocument{\providecommand{\lstlistingname}{}\renewcommand{\lstlistingname}{Listado}%
  \providecommand{\lstlistlistingname}{}\renewcommand{\lstlistlistingname}{Índice de listados}}
\usepackage{amsmath, amssymb}
\usepackage{graphicx}
\usepackage[margin=2.35cm]{geometry}
\usepackage[most]{tcolorbox}
\usepackage{booktabs}
\usepackage{tabularx}
\usepackage{longtable}
\usepackage{array}
\usepackage{float}
\usepackage{xcolor}
\usepackage{hyperref}
\usepackage{fancyhdr}
\setCJKmainfont[AutoFakeBold=2.4]{AR PL UMing CN}
\setCJKsansfont[AutoFakeBold=2.4]{Droid Sans Fallback}
\setCJKmonofont[AutoFakeBold=2.4]{Droid Sans Fallback}
\hypersetup{colorlinks=true, linkcolor=blue!55!black, urlcolor=blue!60!black}
\graphicspath{{./}{figures/}}
\setlength{\parskip}{0.55em}
\setlength{\parindent}{2em}
\setlength{\emergencystretch}{3em}
\renewcommand{\arraystretch}{1.18}
\newtcolorbox{knowledgebox}[1]{enhanced, colback=blue!5!white, colframe=blue!70!black, colbacktitle=blue!70!black, coltitle=white, fonttitle=\bfseries, title={#1}, sharp corners, breakable}
\newtcolorbox{importantbox}[1]{enhanced, colback=yellow!10!white, colframe=yellow!75!black, colbacktitle=yellow!75!black, coltitle=black, fonttitle=\bfseries, title={#1}, sharp corners, breakable}
\newtcolorbox{warningbox}[1]{enhanced, colback=red!5!white, colframe=red!70!black, colbacktitle=red!70!black, coltitle=white, fonttitle=\bfseries, title={#1}, sharp corners, breakable}
\pagestyle{fancy}
\fancyhf{}
\fancyfoot[C]{\thepage}
\fancyfoot[R]{\footnotesize\color{gray}\href{https://github.com/hqhq1025/ai-course-notes}{AI Course Notes}}
\renewcommand{\headrulewidth}{0pt}
\renewcommand{\footrulewidth}{0pt}
\newcommand{\videourl}{https://www.youtube.com/watch?v=khrOsS7YQn4}
\newcommand{\repourl}{https://github.com/hqhq1025/ai-course-notes}

\begin{document}
\begin{titlepage}
\centering
{\Large Notas didácticas de una entrevista a profundidad\par}
\vspace{1.0cm}
{\huge\bfseries Wu Minghui (吴明辉) narra 19 años de historia: empresas ToB, datos privados y Agentic Model empresarial}\par
\vspace{0.5cm}
{\Large Zhang Xiaojun conversa con Wu Minghui, fundador de Mininglamp Technology (明略科技)\par}
\vspace{0.35cm}
{\large Elaborado a partir del video público de Zhang Xiaojun Podcast, la transcripción hecha con GPU, la estructura de capítulos y diagramas conceptuales propios\par}
\vspace{0.3cm}
{\large 2025-10-09\par}
\vspace{0.85cm}
\includegraphics[width=0.88\textwidth,height=0.42\textheight,keepaspectratio]{cover.jpg}\par
\vfill
\begin{tcolorbox}[width=0.92\textwidth, colback=black!2!white, colframe=black!55, sharp corners]
\textbf{Título del video}: 116. Wu Minghui narra 19 años de historia: una larga trayectoria de altibajos, giros dolorosos, Agentic Model empresarial, el juego numérico del mundo real, IPO\par
\textbf{Canal del video}: Zhang Xiaojun Podcast\par
\textbf{Duración del video}: 03:47:45\par
\textbf{Enlace del video}: \href{\videourl}{\nolinkurl{\videourl}}\par
\textbf{Nota sobre los materiales}: el video de YouTube no tiene subtítulos disponibles; el texto principal se elaboró a partir de la transcripción con faster-whisper large-v3 en CUDA, los capítulos del video, la portada y figuras didácticas propias. Las tomas fijas de la entrevista no se reutilizan en el texto principal.\par
\textbf{Página del proyecto}: \href{\repourl}{\nolinkurl{\repourl}}\par
\end{tcolorbox}
\end{titlepage}

\tableofcontents
\newpage
\section{Guía de lectura: por qué la historia de una empresa ToB también es una lección sobre la industria de la IA}

Esta sección explica primero por qué vale la pena incluir este episodio en la serie general de IA/internet. En apariencia, en el EP116 Wu Minghui (吴明辉), antes de la salida a bolsa, repasa Minglue Keji (明略科技), Miaozhen Xitong (秒针系统), las fusiones y adquisiciones, el flujo de efectivo y 19 años de historia de una empresa ToB. En un nivel más profundo, lo que discute son las condiciones reales para implantar la IA empresarial: datos privados, workflow del sector, entregas auditables, capacidades del lado de la oferta, Agentic Model, Tool Use y confianza mutua entre personas y máquinas.

Lo más importante de este episodio no es «cómo una empresa resistió hasta su IPO», sino que ofrece un caso ToB que contradice el relato triunfalista del internet. El internet ToC habla de tráfico, crecimiento y efectos de red; los servicios empresariales ToB hablan del trabajo en las instalaciones del cliente, el acceso a los datos, la calidad de la entrega, el flujo de efectivo, la resistencia de la organización y la credibilidad a largo plazo. Con la llegada de la IA, estos viejos problemas no desaparecen, sino que se reordenan: entrenar modelos base con datos públicos y vender tokens será un terreno de competencia cada vez más feroz, y solo las empresas que poseen datos privados y procesos del sector podrán diferenciarse en el Agentic Model empresarial.

\begin{figure}[H]
\centering
\includegraphics[width=0.96\textwidth]{company-history-map.png}
\caption{Los 19 años de historia de Minglue: el primer emprendimiento, el segundo emprendimiento, el giro doloroso y la IA empresarial vuelven a unirse en una sola línea. Diagrama conceptual de elaboración propia, a partir del contenido de la conversación (00:02:11--03:47:40).}
\end{figure}

\begin{knowledgebox}{Lectura de la figura: no es una historia emprendedora lineal}
La primera etapa corresponde a Miaozhen, los sistemas de recomendación, la medición de publicidad y la complejidad técnica de la Web2; la segunda, a Minglue, el análisis de datos al estilo de Palantir, la ruta ToB/ToG y las fusiones y adquisiciones; el giro doloroso trata del flujo de efectivo, la concentración del negocio y la capacidad de gestión; la era de la IA reinterpreta toda esa acumulación como datos privados, workflow empresarial y Agentic Model.
\end{knowledgebox}

\begin{importantbox}{Tesis central del episodio}
Lo decisivo en la IA empresarial no es «venderle un modelo general a la empresa», sino conectar el modelo con los datos privados, los sistemas de herramientas, los procesos de negocio y los límites de responsabilidad de la empresa, para que la IA produzca, en el trabajo real, resultados auditables, entregables y con un precio que pueda cobrarse.
\end{importantbox}

\begin{knowledgebox}{Ruta de lectura del episodio}
\begin{tabularx}{\textwidth}{p{0.22\textwidth}p{0.32\textwidth}X}
\toprule
Línea & Contenido principal & Pregunta que conviene llevarse \\
\midrule
Historia de la empresa & Miaozhen, Minglue, fusiones y adquisiciones, giro doloroso & ¿Por qué una empresa ToB necesita acumular durante mucho tiempo en lugar de depender de un solo estallido de crecimiento? \\
Línea de datos & Datos de comportamiento de la Web2, conversational data, data privada & ¿Qué datos pueden convertirse en un factor de diferenciación para la IA empresarial? \\
Línea de producto & DeepMiner, Agentic Model, Tool Use & ¿Cómo pasa la IA de ser una herramienta de análisis a ser un sistema que ejecuta flujos de trabajo? \\
Línea de organización & Capacidades del lado de la oferta, dueños/socios, IA confiable & ¿Cómo reescribirá la IA la organización de las empresas y sus límites de responsabilidad? \\
\bottomrule
\end{tabularx}
\end{knowledgebox}

\subsection{Resumen de la sección}

La manera de leer el EP116 es tratar la historia de la empresa como un caso de negocio tecnológico. La primera mitad explica por qué son difíciles el acceso a los datos y la gestión de un negocio ToB; la segunda explica por qué, en la era de la IA, esas mismas dificultades se convierten en barreras de entrada.
\section{Primer emprendimiento: sistemas de recomendación, Miaozhen y la complejidad técnica de Web2}

Esta sección retoma la idea de la introducción de que «la historia de la empresa es también un caso de negocio tecnológico» y responde primero una pregunta: por qué las capacidades de datos de los inicios de Web2 anticiparon la IA empresarial. El primer plan de negocios de Wu Minghui (吴明辉) era un recommendation system, no un sistema publicitario. Este detalle es importante, porque muestra que al equipo le interesaban desde el principio el contenido personalizado, los ID de usuario, la recomendación y los sistemas dinámicos propios de la era Web2. Las páginas de Web1 pueden almacenarse en caché y todos ven el mismo contenido; a partir de Web2, cada usuario ve su propio contenido y la complejidad del sistema aumenta de forma notable.

\begin{importantbox}{Los nuevos problemas que Web2 trajo a los equipos técnicos}
Cuando cada persona ve un contenido distinto, la caché, las bases de datos, la personalización, la recomendación, los registros y el cómputo en tiempo real se vuelven más complejos. Los primeros sistemas de recomendación y de medición publicitaria no eran modelos grandes en el sentido actual, pero ya se ocupaban de «comprender el comportamiento en internet a partir de los datos».
\end{importantbox}

\begin{knowledgebox}{Términos clave: la continuidad de Web2 a la IA empresarial}
\begin{tabularx}{\textwidth}{p{0.22\textwidth}p{0.32\textwidth}X}
\toprule
Concepto & Problema que resolvía entonces & Por qué se relaciona con la IA actual \\
\midrule
User ID & Que el sistema sepa «quién está viendo» & La IA empresarial también debe saber a quién pertenece una tarea, a quién corresponden los permisos y a quién pertenece el contexto. \\
Sistema de recomendación & Asignar contenido distinto a cada usuario & Un Agentic Model también debe generar sugerencias de acción distintas para cada rol. \\
Medición publicitaria & Que varias partes confíen en los mismos datos & La salida de la IA empresarial debe poder auditarse; de lo contrario, no puede entrar en las decisiones de negocio. \\
Registros y datos de comportamiento & Registrar las interacciones reales & El valor de los datos privados proviene de registrar de forma continua los procesos de negocio. \\
\bottomrule
\end{tabularx}
\end{knowledgebox}

Desde esta perspectiva, el primer emprendimiento no es una «precuela anterior a la IA», sino el entrenamiento de base para la IA empresarial. Los sistemas de recomendación desarrollaron la capacidad de distribución personalizada; la medición publicitaria, la capacidad de medición confiable; y la infraestructura de Web2, la capacidad de manejar alta concurrencia y el contexto del usuario. Con el tiempo, estas capacidades cambiaron de nombre: la personalización se convirtió en context, la medición confiable en capacidad de auditoría y los sistemas de alta concurrencia en la entrega estable de workflows empresariales.

\begin{knowledgebox}{La transferencia de capacidades antiguas a capacidades nuevas}
\begin{tabularx}{\textwidth}{p{0.24\textwidth}p{0.32\textwidth}X}
\toprule
Capacidad antigua & Forma de negocio de entonces & Transferencia en la era de la IA \\
\midrule
Recomendación personalizada & Mostrar contenido distinto a cada usuario & Generar sugerencias de tareas distintas para cada rol, departamento y cliente. \\
Medición publicitaria & Medir tráfico, exposición y conversión & Medir si las acciones del Agent producen resultados de negocio reales. \\
Tercero neutral & Que varias partes reconozcan el mismo conjunto de datos & La salida de la IA empresarial debe ser aceptable a la vez para ventas, el área legal, finanzas y los clientes. \\
Formación en subcontratación y entrega & Construir sistemas para clientes como Sina (新浪) y Sohu (搜狐) & Entender las necesidades reales de los clientes, en lugar de limitarse a hacer demos de laboratorio. \\
\bottomrule
\end{tabularx}
\end{knowledgebox}

\begin{knowledgebox}{Los tres hilos ocultos del primer emprendimiento}
\begin{tabularx}{\textwidth}{p{0.20\textwidth}p{0.34\textwidth}X}
\toprule
Hilo oculto & Lo que parecía entonces & Lo que resultó después \\
\midrule
Hilo técnico & Entrega de ingeniería para portales, publicidad y sistemas de clientes & Aprender a manejar datos y confiabilidad dentro de sistemas de negocio reales. \\
Hilo comercial & Paso de la subcontratación y los sistemas de recomendación a la medición publicitaria & Aprender a encontrar una capa intermedia confiable en transacciones entre varias partes. \\
Hilo psicológico & Competencias de matemáticas y la aleatoriedad del emprendimiento & Aprender a tratar la incertidumbre de largo plazo como un problema resoluble. \\
\bottomrule
\end{tabularx}
\end{knowledgebox}

\subsection{El éxito de Miaozhen: de la tecnología a la medición neutral}

Miaozhen Systems (秒针系统) tuvo éxito porque los anunciantes, los medios y las agencias necesitaban una parte neutral que hiciera la medición. El valor de ese papel reside en la confiabilidad de los datos, la uniformidad de los estándares y la neutralidad comercial. Wu Minghui también mencionó que, si Miaozhen se hubiera puesto a comprar y revender tráfico por su cuenta, habría perdido su papel de árbitro neutral. Este es un problema de límites muy típico de las empresas ToB: la tentación de los ingresos a corto plazo puede dañar la confianza a largo plazo.

\begin{knowledgebox}{La lógica comercial de la medición neutral}
Lo que buscan los clientes de una empresa de medición publicitaria no es un algoritmo que «parezca más inteligente», sino una base de hechos que todas las partes puedan reconocer. Una vez que los datos confiables se convierten en el estándar de las transacciones, representan una barrera de entrada mayor que cualquier funcionalidad aislada.
\end{knowledgebox}

\subsection{De las competencias de matemáticas a la psicología del emprendimiento}

En la entrevista se mencionan varias veces el entrenamiento para competencias de matemáticas, las olimpiadas y el ajuste psicológico. Esto explica la resistencia de Wu Minghui como emprendedor: ver una dificultad como un problema por resolver y confiar en que, con más tiempo, se puede resolver. Esta forma de pensar es fundamental para emprender en ToB, porque en ToB rara vez hay un despegue de la noche a la mañana; predominan la entrega en ciclos largos, la confianza de los clientes y la resiliencia organizacional.

\begin{warningbox}{El doble filo de la mentalidad de competencia}
El entrenamiento para competencias desarrolla una gran capacidad para descomponer problemas y para postergar la gratificación, pero también puede llevar a un fundador a subestimar la complejidad del mercado, las ventas, la organización y la estructura de capital. Una empresa ToB no solo resuelve problemas técnicos; también tiene que resolver los problemas de clientes, flujo de efectivo y equipo.
\end{warningbox}

\subsection{Resumen de la sección}

El primer emprendimiento muestra que Wu Minghui no llegó a los servicios empresariales desde el «concepto de IA», sino desde los datos de Web2, la medición publicitaria, la medición neutral y los sistemas de recomendación. Estas capacidades tempranas se transformaron después en la intuición de fondo sobre los datos privados empresariales y el Agentic Model.
\section{Segundo emprendimiento: la visión de Palantir, las adquisiciones y un giro doloroso}

La sección anterior trató sobre Miaozhen (秒针); esta sección entra en el segundo emprendimiento. Minglüe Shuju (明略数据) llamó la atención de Guang Mi (广密), entre otras razones, porque parecía una versión china de Palantir: atendía a clientes gubernamentales y empresariales con integración de datos, plataformas de análisis y entrega de ingeniería. Sin embargo, el mercado ToB chino es distinto del estadounidense: ToG, ToB, la conversión en producto, la entrega personalizada, el flujo de caja y la capacidad organizacional modifican la ruta.

\begin{figure}[H]
\centering
\includegraphics[width=0.94\textwidth]{palantir-to-b-logic.png}
\caption{Palantir y el ToB chino: no se trata de copiar el modelo ToG, sino de encontrar un ciclo cerrado de datos, escenarios y entrega. Diagrama conceptual propio, elaborado a partir de la conversación entre 00:56:08 y 01:15:00.}
\end{figure}

\begin{knowledgebox}{Lectura de la figura: Palantir es un punto de referencia, no la respuesta}
Lo que Palantir enseña está en la integración de datos, las plataformas de análisis y una sólida capacidad de entrega; pero la ruta real de Minglüe tiene que resolver los clientes empresariales chinos, los datos sectoriales, los pagos en ToB, la personalización y los costos organizacionales. Copiar el modelo de negocio no suele bastar: hay que volver a encontrar un ciclo cerrado local.
\end{knowledgebox}

\subsection{La adquisición de Red/Xiaohong: no solo comprar una application, sino una puerta de entrada a los datos}

El apartado anterior mostró que Palantir es solo un punto de referencia; este apartado examina cómo Minglüe buscó su propia puerta de entrada a los datos dentro del ecosistema de Qiye Weixin (企业微信). Un pasaje muy importante de la entrevista es la adquisición del equipo de Xiaohong (小红). Lo que Wu Minghui (吴明辉) valoraba no era una simple application, sino la install base de productos como Weiban (微伴) dentro del ecosistema de Qiye Weixin, así como la gran cantidad de conversational data que generan las conversaciones de ventas, de atención al cliente y con los clientes. Esto coincide justamente con la línea de investigación en Conversational Intelligence que siguió durante su doctorado en ingeniería en la Universidad de Pekín.

\begin{figure}[H]
\centering
\includegraphics[width=0.96\textwidth]{conversational-data-acquisition.png}
\caption{La pista de los Conversational Data: una adquisición no solo compra una aplicación, sino una puerta de entrada a los datos y escenarios de IA. Diagrama conceptual propio, elaborado a partir de la conversación entre 00:04:36 y 00:07:40.}
\end{figure}

\begin{knowledgebox}{Lectura de la figura: por qué los datos de conversación tienen valor}
Las conversaciones de ventas y de atención al cliente de una empresa no son registros de chat comunes, sino la huella real de las necesidades del cliente, el manejo de objeciones, el avance de las transacciones, las explicaciones de producto y la calidad del servicio. Son a la vez datos de entrenamiento y una puerta de entrada al flujo de trabajo.
\end{knowledgebox}

\begin{importantbox}{La lógica de IA de una adquisición}
Si solo se miran los ingresos, adquirir una application puede ser apenas ampliar la línea de productos; vista desde la IA, adquirir una puerta de entrada a un flujo de trabajo de alta frecuencia equivale a obtener un canal que produce datos privados de forma continua. Lo segundo es lo que tiene valor a largo plazo.
\end{importantbox}

\subsection{2020--2022: demasiados frentes abiertos y presión sobre el flujo de caja}

La segunda mitad del segundo emprendimiento fue un giro doloroso. La empresa abrió varios frentes a la vez, la presión sobre el flujo de caja aumentó, el entorno del mercado de capitales cambió y el fundador empezó a pasar de ser un founder idealista a ser un founder que sabe administrar. La lección es muy instructiva: si una empresa ToB tiene una línea de productos demasiado amplia, una entrega demasiado pesada y una organización demasiado dispersa, el flujo de caja y la complejidad de la gestión la castigan.

\begin{figure}[H]
\centering
\includegraphics[width=0.96\textwidth]{two-b-company-turnaround.png}
\caption{El giro doloroso de una empresa ToB: los frentes demasiado amplios, la presión sobre el flujo de caja y la capacidad de gestión obligan a la empresa a volver a enfocarse. Diagrama conceptual propio, elaborado a partir de la conversación entre 01:15:00 y 01:45:01.}
\end{figure}

\begin{knowledgebox}{Lectura de la figura: el giro no es un fracaso, sino una mejora de la capacidad de gestión}
Abrir demasiados frentes genera presión de costos; los años difíciles obligan a la empresa a volver a enfocarse, y solo después de enfocarse puede generar recursos propios con sus productos y servicios. Para una empresa ToB, sobrevivir ya es en sí parte de su capacidad.
\end{knowledgebox}

\subsection{Respeto al capital: una valuación más alta no siempre es mejor}

Después de replegar los frentes, este apartado vuelve a la estructura de capital. Wu Minghui comenta que muchos emprendedores solo llegan a respetar de verdad al capital después de pasar por un ciclo así. Una valuación alta y mucho financiamiento parecen reducir la dilución, pero si la escala del negocio y la ruta hacia la salida a bolsa no los acompañan, aparecen recompras, presión de los inversionistas y dificultades para el fundador. Esto vale en particular para las empresas de software ToB: mientras los ingresos no alcanzan cierta escala, la salida por el mercado de capitales es muy estrecha.

\begin{warningbox}{Una valuación alta no es una ventaja competitiva duradera}
La valuación es una compresión del flujo de caja futuro y de las expectativas del mercado de capitales, no la capacidad real de la empresa. Cuanto más financiamiento, más fácil es que la organización crezca demasiado rápido; cuanto más alta la valuación, mayor la presión futura de recompras y de salida. Si una empresa ToB no puede generar recursos propios, la valuación alta se convierte en una deuda.
\end{warningbox}

\begin{knowledgebox}{Comparación de etapas de la empresa: tres capacidades dentro de una misma compañía}
\begin{tabularx}{\textwidth}{p{0.20\textwidth}p{0.30\textwidth}X}
\toprule
Etapa & Problema principal & Capacidad desarrollada \\
\midrule
Etapa Miaozhen & Monitoreo publicitario, medición neutral, datos de la Web2 & Conciencia de datos confiables y de estándares de transacción. \\
Etapa Minglüe & Análisis de datos ToB/ToG, visión de Palantir, adquisiciones & Capacidad de manejar datos sectoriales, entregar e integrar organizaciones. \\
Etapa de IA & Datos privados, Agentic Model, ejecución confiable & Convertir los datos y el workflow en entregas mediante agentes. \\
\bottomrule
\end{tabularx}
\end{knowledgebox}

\subsection{Resumen de la sección}

El segundo emprendimiento muestra la transición de una empresa tecnológica a una empresa que sabe administrarse. La visión de Palantir dio la dirección y las adquisiciones aportaron una puerta de entrada a los datos, pero el flujo de caja y el exceso de frentes obligaron a la empresa a aprender a enfocarse.
\section{IA de nivel empresarial: los token públicos llevan a una competencia feroz; los datos privados tienen valor}

Las dos secciones anteriores explicaron por qué Minglüe (明略) cuenta con puntos de acceso a datos, experiencia de entrega y resistencia operativa; esta sección entra en el núcleo técnico de todo el episodio, y la pregunta pasa a ser: cómo se convierte esa acumulación, en la era de la IA, en un Agentic Model de nivel empresarial. Wu Minghui (吴明辉) considera que las empresas que entrenan modelos base con datos públicos y cuyo modelo de negocio es vender token enfrentarán una competencia feroz, que fácilmente terminará reduciéndose al costo de la electricidad. Esto se debe a que los datos públicos y las capacidades de los modelos generales tienden a converger, y el precio del token bajará. La diferenciación de la IA empresarial no está en volver a entrenar un modelo público, sino en los datos privados, los procesos de la industria y los sistemas de herramientas reales.

\begin{figure}[H]
\centering
\includegraphics[width=0.94\textwidth]{private-data-value.png}
\caption{Valor diferencial de los datos privados: vender token con datos públicos lleva a una competencia feroz; los datos privados pueden generar valor para la industria. Diagrama conceptual propio, elaborado a partir de la conversación de 01:45:01--02:06:40.}
\end{figure}

\begin{knowledgebox}{Lectura de la figura: por qué los datos privados generan un sobreprecio}
Las capacidades entrenadas con datos públicos se vuelven cada vez más una mercancía; los datos privados de una empresa incluyen registros de clientes, procesos, transacciones, contratos, conversaciones, equipos y riesgos. Solo cuando el modelo entra en ese contexto puede generar resultados diferenciados por los que la empresa esté dispuesta a pagar.
\end{knowledgebox}

\begin{warningbox}{Los «datos privados» no consisten simplemente en subir los archivos de la empresa}
Lo verdaderamente difícil de los datos privados de una empresa está en los permisos, la actualización en tiempo real, la estructuración, los datos sucios, el significado de negocio y los límites de responsabilidad. Meter los documentos en una base de datos vectorial es solo el primer paso; lograr que el Agent entienda correctamente quién puede ver, quién puede modificar y qué resultado es entregable es el verdadero umbral de la IA empresarial.
\end{warningbox}

\begin{knowledgebox}{Arquitectura de cuatro capas de la IA empresarial}
\begin{tabularx}{\textwidth}{p{0.20\textwidth}p{0.32\textwidth}X}
\toprule
Capa & Función & Forma típica de falla \\
\midrule
Capa de datos & Integrar datos privados, permisos e historial & Datos sucios, permisos erróneos, contexto incompleto. \\
Capa de modelo & Comprender, razonar y generar acciones candidatas & Solo sabe resumir, no entiende las restricciones del negocio. \\
Capa de herramientas & Invocar sistemas de negocio, bases de datos e interfaces de procesos & No puede escribir de vuelta en el sistema; se queda en sugerencias. \\
Capa de gobernanza & Auditoría, revisión, responsabilidad y retroalimentación & La empresa no se atreve a delegar y no es posible desplegar a escala. \\
\bottomrule
\end{tabularx}
\end{knowledgebox}

\begin{knowledgebox}{Cadena de conversión del modelo general a la entrega empresarial}
\begin{tabularx}{\textwidth}{p{0.22\textwidth}p{0.32\textwidth}X}
\toprule
Eslabón de conversión & Pregunta que debe responder & Valor de la experiencia tipo Minglüe \\
\midrule
Semántica de la industria & ¿A qué se refiere exactamente la empresa cuando dice «cliente», «riesgo» o «prospecto»? & La larga trayectoria de entrega ToB acumuló semántica de la industria y comprensión del entorno real del cliente. \\
Estado del proceso & ¿La tarea está en preventa, cumplimiento del contrato, cobranza o posventa? & Los datos de CRM, ventas y conversaciones de atención al cliente aportan la trayectoria del proceso. \\
Medición de resultados & ¿La acción mejora de verdad los ingresos, los costos o el riesgo? & La experiencia de Miaozhen (秒针) y en análisis de datos pone énfasis en que todo sea medible y auditable. \\
Adopción organizacional & ¿Quién aprueba, quién revisa, quién asume las consecuencias? & La operación ToB formó la conciencia de permisos, responsabilidad y entrega. \\
\bottomrule
\end{tabularx}
\end{knowledgebox}

Aquí también hay implícito un juicio sobre el modelo de negocio: si una empresa de modelos solo vende token públicos, el precio se verá presionado continuamente a la baja por el aumento de la oferta y la optimización de la inferencia; si una empresa de servicios empresariales convierte los token en acciones de negocio con precio asignable, el precio unitario ya no depende solo del costo de cómputo. Por ejemplo, con el mismo consumo de un millón de token, generar un resumen ordinario y descubrir un riesgo contractual grave tienen un valor comercial completamente distinto. La unidad de precio de la IA empresarial podría terminar pasando del token a las tareas, los riesgos, el aumento de ingresos y el ahorro de costos.

\begin{importantbox}{Del precio por token al precio por resultado}
La unidad de cobro de las API de modelos públicos suele ser el token; la unidad de cobro ideal de la IA empresarial debería ser el resultado. Solo cuando la salida del modelo puede entrar en el workflow y producir beneficios de negocio verificables, la empresa pagará por «resultados» y no por «número de palabras».
\end{importantbox}

\subsection{El juego numérico del mundo real}

El apartado anterior trató los datos privados; este explica por qué es difícil convertir directamente los datos empresariales en capacidades del modelo. Wu Minghui llama a los servicios empresariales el juego numérico del mundo real. El go se presta muy bien a la IA porque el tablero es transparente, las reglas son completas y el context es total; el trabajo empresarial es lo contrario: el contexto está incompleto, la información está dispersa, los objetivos son ambiguos, los procesos cruzan sistemas y los resultados deben responder ante clientes reales y ante las finanzas. La dificultad de la IA empresarial no es «si sabe responder», sino si puede convertir un contexto incompleto en acciones ejecutables.

\begin{importantbox}{Fórmula de la dificultad de la IA empresarial}
\[
\text{Enterprise AI Value} \approx f(\text{Private Data}, \text{Workflow}, \text{Tool Use}, \text{Trust})
\]
Donde \(\text{Private Data}\) es el contexto privado de la empresa, \(\text{Workflow}\) son los procesos de negocio, \(\text{Tool Use}\) es la capacidad de invocar sistemas y herramientas, y \(\text{Trust}\) es la confianza de que todo sea auditable, controlable y entregable.
\end{importantbox}

\begin{knowledgebox}{Términos clave: las cuatro restricciones de la IA empresarial}
\begin{tabularx}{\textwidth}{p{0.20\textwidth}p{0.32\textwidth}X}
\toprule
Restricción & Significado & Consecuencia para el producto \\
\midrule
Private Data & Datos de clientes, procesos, transacciones y conversaciones exclusivos de la empresa & Determina si el modelo entiende el entorno real del negocio. \\
Workflow & Procesos reales que cruzan departamentos y sistemas & Determina si el Agent puede pasar de la sugerencia a la entrega. \\
Tool Use & Invocar herramientas como CRM, ERP, sistemas de tickets y bases de datos & Determina si la IA puede cambiar el estado del negocio. \\
Trust & Permisos, auditoría, revisión y límites de responsabilidad & Determina si la empresa se atreve a confiarle tareas a la IA. \\
\bottomrule
\end{tabularx}
\end{knowledgebox}

Por eso «el juego numérico del mundo real» es más difícil que el go. El espacio de estados del go es enorme, pero el tablero está completo, las reglas son claras y el resultado se puede determinar; el espacio de estados del negocio empresarial no solo es grande, sino que muchas variables están ocultas en la mente de las personas, en las relaciones con los clientes, en los contratos históricos y en los procesos entre sistemas. Para que un Agent pueda trabajar aquí, primero debe convertir el conocimiento tácito en memory aprovechable por el sistema y luego combinar esa memory con la invocación de herramientas.

\subsection{Agentic Model: de las capacidades del modelo a la ejecución de flujos de trabajo}

Este apartado explica el Agentic Model. En la entrevista, Wu Minghui ve al Agent o al Tool Use como el nuevo vínculo en el ámbito de los servicios empresariales. El modelo ya no solo responde preguntas, sino que se conecta a los datos, comprende la tarea, invoca herramientas, completa el proceso y devuelve resultados auditables. Esto se relaciona con el Agent loop de EP139 y el Agent OS de EP118, pero el escenario empresarial pone más énfasis en la implementación privada, los permisos, la auditoría y los límites de responsabilidad.

\begin{figure}[H]
\centering
\includegraphics[width=0.96\textwidth]{enterprise-agentic-model.png}
\caption{Agentic Model de nivel empresarial: de los datos privados a la invocación de herramientas y, después, a la ejecución de los flujos de trabajo de la empresa. Diagrama conceptual propio, elaborado a partir de la conversación de 01:45:01--03:00:00.}
\end{figure}

\begin{knowledgebox}{Lectura de la figura: un Agent empresarial no es un chatbot}
De izquierda a derecha, el Agent empresarial primero debe obtener los datos privados y el contexto de la empresa, luego usar el modelo para comprender y razonar, después invocar mediante Tool Use sistemas como CRM, ERP, sistemas de tickets y bases de datos, y por último entrar en un flujo de trabajo auditable. Si falta cualquier eslabón, es muy difícil que se convierta en un producto de nivel empresarial.
\end{knowledgebox}

\subsection{DeepMiner: extraer de los datos empresariales pistas para la acción}

El Agentic Model resuelve el «cómo ejecutar»; este apartado usa DeepMiner para explicar «dónde se descubren las tareas». DeepMiner aparece poco en la entrevista, pero es una pista para entender la dirección de los nuevos productos de Minglüe. Representa una forma de producto de IA empresarial: no se limita a preguntas y respuestas, sino que integra datos de múltiples sistemas, extrae reglas, riesgos y oportunidades, genera hallazgos y luego ejecuta las acciones siguientes mediante un Agent o herramientas.

\begin{figure}[H]
\centering
\includegraphics[width=0.96\textwidth]{deepminer-product-loop.png}
\caption{Ciclo cerrado del producto DeepMiner: usar la IA para extraer reglas, riesgos y pistas para la acción de los datos de la empresa. Diagrama conceptual propio, elaborado a partir de la conversación de 02:31:58--03:36:59.}
\end{figure}

\begin{knowledgebox}{Lectura de la figura: de la minería de datos a la entrega Agentic}
El BI y el análisis de datos tradicionales se orientan más a mostrar hallazgos; un producto tipo DeepMiner, si se conecta a un Agent, puede ir más allá de «descubrir el problema» y llegar a «proponer acciones», «invocar sistemas» y «registrar resultados». Este es precisamente el punto decisivo para que la IA empresarial pase de ser una herramienta de análisis a ser un sistema de producción.
\end{knowledgebox}

\begin{importantbox}{El cambio decisivo de los productos tipo DeepMiner}
El análisis de datos de antes solía quedarse en dashboard e informes; un producto Agentic de nivel empresarial debe responder además: qué anomalía se descubrió, qué acción debe tomarse, quién aprueba la acción, qué sistema se invoca y cómo regresa el resultado.
\end{importantbox}

\begin{knowledgebox}{De Copilot a Agentic System}
\begin{tabularx}{\textwidth}{p{0.22\textwidth}p{0.34\textwidth}X}
\toprule
Etapa & Forma típica & Requisitos del lado de la empresa \\
\midrule
Copilot & La persona pregunta, la IA responde y la persona copia el resultado & Asistencia de bajo riesgo y alta frecuencia; mejora sobre todo la eficiencia individual. \\
Workflow Agent & La IA lee datos, llena formularios e inicia procesos & Debe integrar permisos, herramientas y aprobaciones. \\
Agentic System & Varios Agent colaboran para completar tareas & Requiere registros, orquestación, recuperación ante fallas y responsabilidad sobre los resultados. \\
Trusted Agentic Model & Asume de forma confiable tareas críticas de la empresa & Debe ser auditable, revisable y desplegable de forma privada o ejecutarse de forma controlada. \\
\bottomrule
\end{tabularx}
\end{knowledgebox}

\subsection{Resumen de la sección}

El núcleo de la IA empresarial está en los datos privados y los flujos de trabajo. Los token públicos se abaratarán; lo que realmente tiene un sobreprecio es el Agentic Model capaz de convertir el contexto de la empresa en resultados.
\section{Nuevas relaciones de producción: capacidad del lado de la oferta, redes de conexión y reescritura de la organización}

La sección anterior trató sobre modelos y productos; esta trata sobre las relaciones de producción. Wu Minghui (吴明辉) sostiene que la IA reestructurará la industria del software, los servicios empresariales y el internet industrial, y que generará nuevas relaciones de conexión y de producción. Las más expuestas son las empresas que solo ofrecen una red de conexión, no tienen capacidad del lado de la oferta y cuya red, además, no es un nodo raíz. La razón es que la IA reducirá el costo de conectar y de procesar información, lo que debilitará el valor de la capa intermedia.

\begin{figure}[H]
\centering
\includegraphics[width=0.94\textwidth]{supply-demand-relation.png}
\caption{Capacidad del lado de la oferta frente a red de conexión: una red que solo conecta, sin capacidad del lado de la oferta, queda en una posición de riesgo. Diagrama conceptual propio, elaborado a partir de la conversación de 02:31:36--02:34:28.}
\end{figure}

\begin{knowledgebox}{Lectura de la figura: la IA le pondrá un nuevo precio a la capa intermedia}
Si una empresa solo se dedica a la intermediación, a la distribución o a mantener cadenas de relaciones, pero no tiene capacidad de producto, datos, herramientas, modelos ni entrega, su valor se reducirá en cuanto un Agent pueda conectar directamente la demanda con la oferta. La capacidad del lado de la oferta es la capacidad de producir resultados reales.
\end{knowledgebox}

\subsection{Las nuevas relaciones de producción: personas, modelos, herramientas y datos se reparten el trabajo de otra forma}

La sección anterior explicó cómo se le pone un nuevo precio a la capa intermedia; esta examina cómo se descompone el trabajo mismo. Las nuevas relaciones de producción que trae la IA no consisten en que «la IA sustituya a todos», sino en una nueva división del trabajo entre personas, modelos, herramientas, datos y organizaciones. El modelo se encarga de la cognición, las herramientas aportan las interfaces de ejecución, los datos proporcionan el contexto, el Agent se ocupa de la colaboración automática y las personas asumen los objetivos, el juicio y la responsabilidad. El futuro de los servicios empresariales no se limita a las API: consiste en reorganizar la producción en torno al resultado de cada tarea.

\begin{figure}[H]
\centering
\includegraphics[width=0.92\textwidth]{new-production-relation.png}
\caption{Las nuevas relaciones de producción que trae la IA: modelos, herramientas, Agent, datos privados y socios humanos se reparten el trabajo de otra forma. Diagrama conceptual propio, elaborado a partir de la conversación de 00:00:42--00:00:55 y 02:31:36--02:34:28.}
\end{figure}

\begin{importantbox}{Criterio para reconocer un cambio en las relaciones de producción}
Si la IA solo mejora la eficiencia de un punto aislado, se trata de una mejora de herramientas; si la IA cambia quién posee los datos, quién ejecuta las tareas, quién asume la responsabilidad y quién obtiene los beneficios, se trata de un cambio en las relaciones de producción.
\end{importantbox}

\begin{knowledgebox}{Términos clave: conexión, lado de la oferta y relaciones de producción}
\begin{tabularx}{\textwidth}{p{0.22\textwidth}p{0.32\textwidth}X}
\toprule
Término & Significado en la entrevista & Cambio en la era de la IA \\
\midrule
Red de conexión & Red que conecta clientes, comercios, información o servicios & El Agent reduce el costo de conectar y el valor de la capa intermedia se reduce. \\
Capacidad del lado de la oferta & Capacidad de producir efectivamente bienes, servicios, modelos o resultados entregables & Cuanto más cerca del resultado, más difícil es prescindir de ella. \\
Relaciones de producción & Quién organiza la producción, quién la ejecuta y quién reparte los beneficios y la responsabilidad & Con la incorporación de modelos y Agent, los roles se redistribuyen. \\
Socio & Persona capaz de responder por los resultados, definir problemas y coordinarse con Agent & Disminuyen los puestos de ejecución repetitiva y aumentan los roles de alta responsabilidad. \\
\bottomrule
\end{tabularx}
\end{knowledgebox}

\subsection{Dueños y socios: una visión de la forma organizacional}

En la entrevista aparece también un juicio extremo: en el futuro, las empresas podrían tener solo dos tipos de personas, dueños y socios. Aquí «socio» no designa a un empleado tradicional, sino a alguien que responde por los resultados junto con la IA, los Agent y las herramientas. Puede que este juicio no se cumpla al pie de la letra, pero expresa una tendencia: cuando la ejecución quede automatizada por la IA, la organización dará más peso a la definición de objetivos, la integración de recursos, la asunción de responsabilidades y la colaboración de alta densidad.

\begin{figure}[H]
\centering
\includegraphics[width=0.94\textwidth]{boss-partner-org.png}
\caption{Organización de dueños y socios: la empresa podría pasar de una organización de empleados a una de pocos dueños y muchos socios de IA. Diagrama conceptual propio, elaborado a partir de la conversación de 02:53:20--03:16:40.}
\end{figure}

\begin{warningbox}{«Solo dueños y socios» no significa que la organización desaparezca de inmediato}
Las empresas reales siguen necesitando puestos, procesos, cumplimiento normativo y gestión. Sin embargo, a largo plazo, los puestos de ejecución repetitiva se reducirán por efecto de la IA, y las personas capaces de definir problemas, integrar recursos y responder por los resultados se parecerán cada vez más a socios.
\end{warningbox}

\subsection{Resumen de la sección}

El impacto de la IA en los servicios empresariales no se limita a una mejora de eficiencia. Reescribirá las redes de conexión, la capacidad del lado de la oferta, la estructura organizacional y el reparto de responsabilidades. Las empresas realmente expuestas son las que solo aportan valor como capa intermedia y carecen de capacidad productiva.
\section{Misión e IA confiable: de inteligente a confiable}

Esta sección concluye con el objetivo personal de Wu Minghui (吴明辉). Cuenta que, antes de la aparición de GPT, su meta de vida era construir una IA más inteligente que él mismo; después se dio cuenta de que la «inteligencia» ya la estaban impulsando las grandes empresas de modelos, así que reorientó su objetivo hacia construir la IA más confiable, es decir, un trusted agentic model. Este cambio de rumbo es decisivo: en la IA empresarial lo esencial no es quién conversa mejor, sino a quién se le puede confiar el uso en operaciones reales del negocio.

\begin{figure}[H]
\centering
\includegraphics[width=0.96\textwidth]{happy-founder-mission.png}
\caption{Alineación de misiones del founder feliz: la misión personal, la misión familiar y la misión de la empresa están estrechamente relacionadas. Diagrama conceptual propio, elaborado a partir de la conversación entre 02:53:20 y 03:47:40.}
\end{figure}

\begin{knowledgebox}{Lectura de la figura: la alineación de misiones es la fuente de energía de un compromiso de largo plazo}
La misión personal determina qué problemas está dispuesto a resolver el fundador, la misión familiar aporta una energía estable y la misión de la empresa se conecta con el valor para el cliente. Wu Minghui concibe las dificultades del emprendimiento como problemas de matemáticas; detrás de esa actitud hay una fuerte capacidad de postergar la gratificación y una alineación de misiones.
\end{knowledgebox}

\subsection{De la más inteligente a la más confiable}

Las secciones anteriores trataron la organización y las relaciones de producción; esta vuelve al criterio último con el que se evalúa la IA empresarial. La competencia entre modelos ToC suele girar en torno a «quién es más inteligente»; los servicios empresariales ToB se preocupan más por «quién es más confiable». La confiabilidad abarca que el origen de los datos sea confiable, que los permisos sean confiables, que el proceso pueda auditarse, que los resultados puedan verificarse, que los errores puedan atribuirse a un responsable y que el sistema admita una implementación privada o controlada. Los clientes empresariales no pagan por un despliegue de virtuosismo técnico; al final pagan por resultados con riesgo controlado.

\begin{importantbox}{Qué significa Trusted Agentic Model}
Un Trusted Agentic Model es un modelo de agente capaz de ejecutar tareas en el contexto privado de una empresa, invocar herramientas, conservar un registro de auditoría, respetar los límites de permisos y entregar resultados verificables. Su objetivo no es dar la respuesta más deslumbrante en una sola ocasión, sino ser el más confiable a lo largo de una operación prolongada.
\end{importantbox}

\begin{knowledgebox}{Lista de verificación para aceptar un Agent confiable}
\begin{tabularx}{\textwidth}{p{0.20\textwidth}p{0.32\textwidth}X}
\toprule
Dimensión & Pregunta de verificación & Por qué importa \\
\midrule
Permisos & ¿Solo lee y modifica datos dentro del alcance autorizado? & Evita accesos no autorizados e incidentes de cumplimiento normativo. \\
Auditoría & ¿Cada acción tiene registros y evidencia que permita reconstruirla? & La empresa necesita atribuir responsabilidades y mejorar. \\
Verificación & ¿Las acciones de alto riesgo requieren confirmación humana? & Reduce los errores irreversibles. \\
Retroalimentación & ¿Los resultados vuelven al modelo y al proceso? & Permite que el sistema mejore de forma continua. \\
\bottomrule
\end{tabularx}
\end{knowledgebox}

\subsection{El empresario feliz: tratar el suffering como un problema por resolver}

La sección anterior abordó el objetivo de producto de la IA confiable; esta vuelve a la persona que sostiene ese objetivo de largo plazo. Al cierre, Wu Minghui dice que es un founder feliz: aunque 2022 fue muy difícil, no tuvo un suffering especial, porque estaba convencido de que sin duda lo resolvería; solo que el problema tardaría más en resolverse. No se trata de que fuera fácil, sino de una actitud que abstrae las dificultades de largo plazo como problemas por resolver. Esto explica por qué un emprendimiento ToB puede resistir 19 años: sin esa capacidad de postergar la gratificación y sin esa estructura psicológica, es muy difícil atravesar entregas prolongadas, vaivenes del capital y reestructuraciones organizacionales.

\begin{knowledgebox}{La estructura psicológica de un founder ToB}
El emprendimiento ToB rara vez ofrece retroalimentación explosiva a corto plazo; se basa más bien en relaciones de largo plazo con los clientes, restricciones de flujo de efectivo, estabilidad organizacional y credibilidad del producto. Que el founder pueda ver las dificultades como problemas descomponibles, y no como un fracaso de su identidad, determina si podrá seguir iterando.
\end{knowledgebox}

Esto también explica por qué el final de la entrevista no es solo un balance de negocio, sino un balance de misiones. Que la misión personal, la familiar y la de la empresa estén estrechamente relacionadas significa que el fundador no trata a la empresa como un proyecto de corto plazo, sino como parte de su propio problema de largo plazo. Esto es especialmente cierto en las empresas de IA empresarial: la confianza de los clientes, la seguridad de los datos, la entrega en cada industria y los Agent confiables exigen una acumulación de largo ciclo, que no puede completarse aprovechando una sola ventana de capacidad de los modelos.

\begin{importantbox}{El costo real de la visión de largo plazo}
La visión de largo plazo no es un eslogan, sino la disposición a soportar retroalimentación lenta, poca visibilidad externa, presión sobre el flujo de efectivo y ajustes organizacionales repetidos. Si una empresa ToB logra entrar en la era de la IA, por lo general no lo debe a un momento de explosión, sino a los clientes, los datos, la capacidad de entrega y la estructura psicológica del founder acumulados durante más de una década.
\end{importantbox}

\subsection{Resumen de la sección}

El objetivo de Wu Minghui para la IA pasó de «más inteligente» a «más confiable», lo que corresponde justamente a la necesidad central de la IA empresarial. Los clientes empresariales no necesitan solo capacidad de modelo, sino ejecución confiable y entrega sostenida a largo plazo.
\section{Síntesis y ampliación}

Esta sección condensa el EP116 en tres niveles. Primero, el nivel de la historia de la empresa: a lo largo del sistema de recomendación, Miaozhen (秒针), Minglüe Shuju (明略数据), las fusiones y adquisiciones y los giros dolorosos, Minglüe (明略) acumuló capacidades de datos ToB, de entrega y de gestión. Segundo, el nivel técnico y comercial: la diferenciación de la IA empresarial proviene de los datos privados, los flujos de trabajo, el Tool Use, el Agentic Model y la entrega confiable, no de la guerra de precios de los token públicos. Tercero, el nivel organizacional: la IA reescribirá las capacidades del lado de la oferta, las redes de conexión, la estructura de los puestos y la forma en que el founder expresa su misión.

\begin{importantbox}{Ubicar el EP116 en la serie de IA de Zhang Xiaojun (张小珺)}
El EP117 trata de artículos académicos e historia de la tecnología, el EP118 trata de VLA, Agent OS y plataformas, y el EP116 trata de la implementación de la IA en escenarios de servicios empresariales: nos recuerda que el valor industrial de la IA no reside solo en las empresas de modelos, sino también en las empresas ToB que poseen datos privados, procesos sectoriales y responsabilidad de entrega.
\end{importantbox}

\subsection{Takeaways clave}

\begin{enumerate}
\item La ventaja competitiva de la IA empresarial no es el modelo público, sino los datos privados, los procesos y la entrega confiable.
\item En ToB, el Agentic Model debe conectarse con herramientas, permisos, auditoría y workflows reales.
\item Las empresas que solo ofrecen redes de conexión, sin dominar capacidades del lado de la oferta, serán revaloradas por la IA.
\item La dificultad central de emprender en ToB es el flujo de efectivo, la entrega, el enfoque y la resistencia organizacional, no una tecnología aislada.
\item El Trusted Agentic Model se acerca más que «el modelo más inteligente» a la razón real por la que compran los clientes empresariales.
\end{enumerate}

\subsection{Lecturas adicionales}

\begin{itemize}
\item Si le interesa el Agentic Model, puede contrastarlo con la revisión sobre Agent del EP139 y la entrevista a Luo Fuli (罗福莉) del EP138, para entender la transición del Agent desde un loop técnico hacia un sistema organizacional y de posentrenamiento.
\item Si le interesan los datos privados empresariales, puede contrastarlo con la revisión sobre datos del EP134, para entender las diferencias entre la fijación de precios de los datos, la Recipe y los datos de despliegues reales.
\item Si le interesan las plataformas de IA y los sistemas operativos, puede contrastarlo con la entrevista a Li Xiang (李想) del EP118 y con el EP127 AI War, para observar la relación entre el Agent OS, el workflow empresarial y los puntos de entrada de las plataformas.
\end{itemize}

\end{document}
